Puja dan puji syukur ke hadirat Allah Subhaanahu wa ta’aala atas limpahan
rahmat, karunia, serta petunjuk-Nya sehingga tugas akhir berupa penyusunan skripsi
dengan judul "Rancang Bangun Rupiah Digital Ritel Berbasis Hyperledger Fabric" ini telah terselesaikan dengan baik. Dalam hal
penyusunan tugas akhir ini, penulis telah banyak mendapatkan arahan, bantuan, serta
dukungan dari berbagai pihak.
Oleh karena itu, pada kesempatan ini penulis
mengucapkan terima kasih kepada:

\begin{enumerate}
	\item Bapak Prof. Ir. Hanung Adi Nugroho, Ph.D., IPM, SMIEEE., selaku Ketua
	      Departemen Teknik Elektro dan Teknologi Informasi Fakultas Teknik Universitas
	      Gadjah Mada.
	\item Bapak Dr. Ir. Ahmad Nasikun, S.T., M.Sc., IPM., selaku Ketua Program Sarjana
	      Program Studi Teknologi Informasi Fakultas Teknik Universitas Gadjah Mada.
	\item Bapak Teguh Bharata Adji, S.T., M.T., M.Eng., Ph.D., selaku Dosen Pembimbing I
	      yang senantiasa memberikan bimbingan dan masukan yang terbaik dalam
	      pelaksanaan serta penulisan laporan tugas akhir.
	\item Bapak Dr. Ir. Guntur Dharma Putra, S.T., M.Sc., selaku Dosen Pembimbing II
	      yang telah memberikan arahan dan masukan teknis yang terbaik dalam pelaksanaan
	      tugas akhir.
	\item Ibu Dr. Eng. Silmi Fauziati, S.T., M.T., selaku Dosen Pembimbing Akademik yang
	      senantiasa membantu penulis perihal dinamika perkuliahan.
	\item Seluruh dosen dan tenaga pendidik Departemen Teknik Elektro dan Teknologi Informasi
	      Fakultas Teknik Universitas Gadjah Mada yang telah memberikan bekal
	      ilmu kepada penulis selama menjalani masa perkuliahan.
	\item Bapak Sujono, S.E., Ibu Mulyati, S.E., dr. Puji Arifianti Ramadhany, dan
	      Muhammad Feizal Syafaat, S.T., selaku keluarga penulis yang senantiasa
	      memanjatkan doa kepada penulis serta memberikan dukungan, semangat, dan motivasi
	      kepada penulis dalam pengerjaan tugas akhir.
	\item UKM UGM eSports 2024 yang telah memberikan penulis pandangan baru terkait dinamika
	      eSports, dinamika UGM, hingga dinamika negara.
	\item Cornelius Arden Satwika Hermawan, S.T., dan Mahsa Quereda Bahjah, S.T., selaku
	      rekan kelompok mata kuliah Pemrograman Berorientasi Obyek yang terus berlanjut
	      hingga Teknisiade, PPK Ormawa, Bengkel Inovasi, Kerja Praktik. Semoga kita
	      dapat menjadi rekan kerja lagi di kesempatan lain.
	\item WS yang telah menemani selama 24 jam dan membuat hidup penulis menjadi warna
	      warni. Terimakasih sudah membuat kehidupan kuliah menjadi seperti tujuh belas
	      lagi seperti yang dideskripsikan oleh Tulus, "Muda jiwa selamanya muda, Kisah
	      kita abadi selamanya".
	\item Abid, Ale, Caca, Ervina, Lintang, Nabila, Niko, Sanie, Vivian, Yajna, Abang
	      Efrem, Bapa Poli, Abang Rinto, Abang Rino, Mama Fani, dan semuanya yang
	      tidak bisa penulis sebutkan. Terimakasih sudah menjadi keluarga untuk penulis di
	      Desa Iso, Kecamatan Kei Kecil Timur. Terimakasih juga sudah merayakan ulang
	      tahun penulis dengan cara yang sangat berkesan.
	\item Tim Jelajah Kei Kecil 2025, khususnya untuk Almarhum Bagas Amar Hakiki.
	      Terimakasih sudah menjadi sahabat yang selalu menemani penulis bahkan saat
	      berada di fase terendah.
	\item Teman-teman Program Studi Teknologi Informasi Departemen Teknik Elektro dan
	      Teknologi Informasi Universitas Gadjah Mada Angkatan 2022 yang telah menemani
	      dan membantu penulis selama menjalani masa perkuliahan.
\end{enumerate}

Tiada gading yang tak retak. Penulis menyadari bahwa skripsi ini masih jauh dari
kata sempurna. Oleh karena itu, penulis sangat mengharapkan segala bentuk masukan,
saran, dan kritik yang membangun. Akhir kata, penulis berharap semoga tulisan ini dapat
memberikan manfaat dan wawasan tambahan bagi kita semua, aamiin

\begin{flushright}
	\begin{tabular}{c}
		Yogyakarta, Juni 2026 \\
		\vspace{1cm}          \\
		Muhammad Fachrizal Giffari
	\end{tabular}
\end{flushright}
